\cotefile{79}
\scanpage{49}

\unsure{Â|À} vrai dire, pour définir les structures triangulées $\Sigma(i)$ $(i\in\unsure{\Pi|\pi})$, il suffirait de connaître (118) sur l'ens

\[
\underset{\substack{\cap\\ S\times S}}{\widetilde A}
=\{(\deleted{\unsure{x}}s,\deleted{\unsure{y}}t)\in S\times S\mid \deleted{\unsure{x}}s\mathbin{\wedge}\deleted{\unsure{y}}t\in\Pi\}
\tag*{(124)}
\]

en tant qu'application

\[
\widetilde A\longrightarrow\Pi\qquad\deleted{\text{\unsure{applicat|appliquée}}}
\tag*{(125)}
\]

d'une partie \inserted{$\widetilde A$} de $S\times S$ dans $\Pi$, \deleted{l'image} \inserted{telle} \deleted{\unsure{inverse}} \inserted{partie} ne rencontre pas la diagonale et est \deleted{stable} stable par la symétrie $(s,t)\mapsto(t,s)$, donc sa donnée équivaut à celle d'une partie $A$ de $\unsure{\mathscr P|\mathcal P|\mathfrak P}_2(S)$, et la donnée de (125), \inserted{application symétrique,} à celle d'une application

\[
 A\xrightarrow{t}\Pi\qquad t=\text{``type''}
 \tag*{(126)}
\]

dont les fibres sont donc les $A_{\Sigma(i)}$, et on \unsure{aura|a}

\[
 A=\bigcup_{i\in\Pi}\bigl(A_{\Sigma(i)}={\unsure{\wedge|t}}^{-1}(\{i\})\bigr)
 \tag*{(127)}
\]

On peut dire que la donnée d'une famille, indexée par $\Pi$, de structures triangulées sur $S$, déterminées par les structures de graphes correspondantes via (121), équivaut à celle d'une partie $A$ de $\mathscr P_2(S)$ (\unsure{ie.|i.e.} d'une partie $\widetilde A$ de $S\times S$ ne rencontrant pas la diag. et stable par symétrie) et d'une application (surjective) (127), de façon que les fibres $A_i=A_{\Sigma(i)}$

% archive sheet number: 48
\scanpage{50}
\mspage{24}

satisfassent (122) et (123). Dans ces données, la structure de \inserted{$Z'_n$-}torseur sur $\unsure{\Pi|\pi}$ n'est pas encore intervenue.

\deleted{Par contre, quand on fait} \inserted{Quand on fait} intervenir cette dernière, $\Pi$ se plonge canoniquement dans un $L'$, et la question se pose de décrire, \uline{en termes des données précédentes}, une application symétrique (118) qui prolonge (125). La question \uline{est liée} à celle de décrire une opération de $\deleted{\unsure{\Lambda|A}}$ $Z'_n$ sur $S$

% slanted

\margin{pas tellement, finalement !}

\[
 \lambda\longmapsto T_\lambda:\quad Z'_n\mathrel{\unsure{\relbar\joinrel\relbar\mkern-2mu\shortmid|\longrightarrow}}\unsure{\mathfrak S|\mathscr S}_S
 \tag*{(124)}
\]

qui satisfasse

\[
 s\wedge T_\lambda t=\lambda(s\wedge t)\qquad\text{pour }s,t\in S,
 \tag*{(125)}
\]

donc \deleted{a fortiori} en particulier

\[
 \begin{aligned}
 s\wedge T_\lambda t&=\lambda(s\wedge t)\qquad\text{pour }(s,t)\in\widetilde A\\
 \text{i.e.}\quad \unsure{t|\tau}(\{s,T_\lambda t\})&=\lambda\cdot t(\{s,t\})\qquad\text{pour }\{s,t\}\in A.
 \end{aligned}
 \tag*{(125$'$)}
\]

On a vu tantôt qu'une telle \deleted{application} \inserted{opération} \deleted{\unsure{est}} \deleted{\unsure{unique|uniquement}} \inserted{de} $Z'_n$ sur $S$ est uniquement déterminée, du fait que l'on a

\[
 \left\{\begin{aligned}
 &\text{si }t,t'\in S\text{ sont tels que}\\
 &\unsure{\text{Ét}|\text{Et}}_{\Sigma(i)}(t)=\text{Ét}_{\Sigma(i)}(t')\qquad
 \text{pour \deleted{tout} }i\in\Pi\\
 &\text{alors }t=t'\qquad\qquad\text{(ou seulement pour \uline{un} }i\in\Pi\text{)}
 \end{aligned}\right.
 \tag*{(126)}
\]

% archive sheet number: 49
\scanpage{51}

La relation (125$'$) \unsure{montre} qu'on a

% slanted

\margin{\uline{NB} Il revient au $=$ de se donner l'op. de $Z'_n$ sur $\unsure{\Pi|\pi}$ (en faisant un $Z'_n$-torseur), ou l'opération $\lambda\mapsto T_\lambda$ de $Z'_n$ sur $S$, reliées par (125$'$) -- \unsure{chacune|une|l'une} détermine l'autre.}

\[
\begin{aligned}
 \{s,t\}\in A&\Longrightarrow
 \overbrace{\{T_\lambda s,T_\lambda t\}}^{\inserted{\{s',t'\}}}\in A\\
 &\text{et}\qquad \unsure{\tau|t}(\{s',t'\})=\lambda^2\,\tau(\{s,t\}),
\end{aligned}
\tag*{(127)}
\]

donc $T_\lambda$ est un \unsure{autom} de la structure définie par \deleted{\unsure{$(A,Z'_n),\tau$}} \inserted{$\{S,\Pi\text{ tors. sous }Z'_n,A\subset S\times S,\tau\}$,} ou ce qui revient au $=$, de la structure ($S$, $\Pi$ tors. sous $Z'_n$, $\Pi\xrightarrow{\Sigma}$ ens des structures triangulées sur $S$ comme ens de sommets), canoniquement associée à ladite structure.

\deleted{Revenons à la question de décrire $\wedge$ (118) en termes des autres structures, on notera pour ceci que l'on a, si}

% slanted

\margin{il faut y revenir ultérieurement !}

\[
 \deleted{\left\{\begin{smallmatrix}\{s,t,u\}\in\unsure{F|\mathcal F}=\coprod_{i\in\Pi}F_{\Sigma(i)}\\v\in S\end{smallmatrix}\right.}
 \quad\deleted{\Longrightarrow s\wedge v+t\wedge v+u\wedge v-}
 \tag*{\deleted{(128)}}
\]

J'aimerais comprendre comment \inserted{(éventuellement)} \uuline{une des structures} $\Sigma(i)$ \uuline{détermine toutes les autres}, et l'opération de $Z'_n$ sur l'ens de ces structures. \deleted{\unsure{L}}

Soit donc $P_0=(S,A_0,F_0)$ une structure de 2-polyèdre \unsure{connexe.|combin.} régulier, triangulé en l'occurrence, fini par dessus le marché !

Soit

\[
 G_0=\operatorname{Aut}(P_0)\lhook\joinrel\longrightarrow\unsure{\mathfrak S|\mathscr S}_S,
 \tag*{(128)}
\]

% archive sheet number: 50
\scanpage{52}
\mspage{25}

et considérons son \unsure{commutant} dans $\unsure{\mathfrak S|\mathscr S}_S$, soit

\[
 Z=\{\lambda\in\mathfrak S_S\mid\lambda u=u\lambda\quad\forall u\in G_0\}
 \tag*{(129)}
\]

Dans le cas où $P_0=SP_{(M,\{+\unsure{e|\ell}_M,-e_M\})}$, on a donc $G_0\simeq{\mathrm{Gl}}^{\prime}_\Lambda(M)^{\pm}$ opérant sur $S=M^{\unsure{*|+}}/\{\pm\mathrm{id}_M\}$, son commutant sauf erreur est réduit à $Z'_n$ (si $\Lambda=\mathbb Z/n\mathbb Z$).

% slanted

\margin{\uuline{à vérifier}}

\deleted{Soit $\lambda\in Z$, alors} On a

\[
 \deleted{\lambda.A_0}\ \deleted{\unsure{\mathfrak z\subset Z}}
 \quad Z\cap G_0=\operatorname{Cent}(G_0)\overset{\mathrm{def}}{=}
 \unsure{\mathfrak z|z}\subset\operatorname{Cent}(Z)
 \quad\deleted{=\operatorname{Cent}\,G}
 \tag*{(130)}
\]

($\mathfrak z=\deleted{\unsure{\{1\}}}{}_{2}Z'$ dans le cas qui nous occupe)

% slanted

\margin{donc $\mathfrak z=$}

\deleted{donc} introduisons aussi

\[
 Z'=Z/\mathfrak z \qquad \unsure{(}
 \tag*{(131)}
\]

(on a $Z'=Z$ dans le cas qui nous occupe). Comme la structure triangulée \inserted{$\Sigma_0$} sur $S$ est connue quand on connaît $G_0\subset\mathfrak S_S$ et \uline{une} de\deleted{\unsure{s}} \inserted{ses} arêtes, ou \uline{une} de ses faces (puisque $G_0$ est transitif sur $A_0$ et sur $F_0$), on doit avoir, pour $\lambda\in Z'$

\[
 \deleted{\unsure{\{}} A_{\lambda\cdot\Sigma_0}\cap A_{\Sigma_0}=\emptyset,\qquad
 F_{\lambda\Sigma_0}\cap F_{\Sigma_0}=\emptyset
 \tag*{(132)}
\]

% archive sheet number: 51
\scanpage{53}

plus généralement, \deleted{pour $\lambda,\lambda'$}

\[
 \lambda,\lambda'\in Z',\quad\lambda\ne\lambda'\quad\Longrightarrow
 \left\{\begin{aligned}
 A_{\lambda\Sigma_0}\cap A_{\lambda'\Sigma_0}&=\emptyset\\
 F_{\lambda\Sigma_0}\cap F_{\lambda'\Sigma_0}&=\emptyset
 \end{aligned}\right. .
 \tag*{(132bis)}
\]

\deleted{\unsure{On|Donc}} \deleted{plus général} \inserted{D'autre part}, si $\lambda\Sigma_0=\Sigma_\lambda$, de sorte que

\[
 \begin{aligned}
 \lambda&\longmapsto\deleted{\unsure{\Sigma_\lambda}}\ T_\lambda(\Sigma_0)\overset{\mathrm{def}}{=}\Sigma_\lambda\\
 Z'&\longrightarrow\text{ens. des structures triangulées}\\
 &\hspace{3em}\text{régulières sur }S
 \end{aligned}
 \tag*{(133)}
\]

% slanted

\margin{c'est une notation choisie pour $\Sigma_\lambda$ -- elle n'est pas cohérente avec la $\Sigma(i)$ de tantôt \unsure{--}, on a \deleted{\unsure{$Z\lambda$}} $\Sigma_\lambda=\Sigma_{\unsure{s(\lambda^2)}}$}

est une application injective, on a

\[
 \left\{\begin{aligned}
 \operatorname{Aut}(S,\Sigma_\lambda)&=G\qquad\text{pour tout }\lambda\in Z'\\
 \operatorname{Aut}^{+}(S,\Sigma_\lambda)&=G^+
 \end{aligned}\right. .
 \tag*{(134)}
\]

Ceci posé, on dira que deux structures 2-polyédrales \unsure{trianglées|triangulées} régulières $\Sigma,\Sigma'$ sur un ens. de sommet $S$ sont \inserted{(\uline{strictement})} \uline{équivalentes}, si $\exists\lambda\in\unsure{\mathfrak S|\mathscr S}_S$, commutant à $\operatorname{Aut}(S,\Sigma)=G_\Sigma$, telle que $\Sigma'=\lambda.(\Sigma)$ -- ce qui implique que

\[
 \left\{\begin{aligned}
 G_\Sigma&=G_{\Sigma'},\qquad \deleted{\unsure{\Sigma}}\lambda^{-1}\in\unsure{\operatorname{Comm}|\operatorname{comm}|\operatorname{Commut}}_{\mathfrak S_S}(G_{\Sigma'})\\
 &\hspace{5em}\Sigma^{\unsure{\prime\prime|\prime}}=\lambda^{-1}.(\Sigma')
 \end{aligned}\right.
 \tag*{(135)}
\]

% archive sheet number: 52
\scanpage{54}
\mspage{26}

de sorte que la relation envisagée est une relation d'équivalence. Une classe d'équivalence \inserted{$\Pi$ pour cette relation} est alors associée à un groupe d'automorphismes commun $G$ pour les structures envisagées, et se trouve alors \unsure{muni|munie} de façon naturelle d'une structure de torseur sous le groupe

\[
 \left\{\begin{aligned}
 Z'&=Z/\unsure{\mathfrak z|z}\deleted{\unsure{\cap G}},\quad
 Z^{\deleted{\unsure{\prime|4}}}=\operatorname{Comm}_{\unsure{\mathfrak S|\mathscr S}_S}(G)\\
 &&\mathfrak z=Z'\cap G=\operatorname{Cent}(G).
 \end{aligned}\right.
 \tag*{(136)}
\]

Le groupe des automorphismes de $(S,\Pi)$ \deleted{est alors} \inserted{qui induisent l'identité sur $Z'$} s'insère alors dans une suite exacte

\diagram{\begin{tikzcd}[ampersand replacement=\&,column sep=1.35em,row sep=1.5em]
{(137)} \& 1\arrow[r] \& G\arrow[r] \& {\operatorname{Aut}(S,\Pi)}\arrow[r] \& Z'\arrow[r] \& 1\\
 \& \& \mathfrak z\arrow[u,hook]\arrow[r,arrows={Hooks[harpoon]-}] \& Z\arrow[u,hook]\arrow[ur,"{\text{épi can.}}"'] \& \&
\end{tikzcd}}

Mais cette analyse, dans le cas qui nous occupe, ne nous fournit pas, à partir d'un $SP_{(M,\{\pm\unsure{e|\ell}_M\})}=P_{\Sigma_0}$, \deleted{toutes} \uline{les} structures $SP_{M,i}$ $i\in{Z'_n}$, mais seulement les structures pour $i\in(Z'_n)^2=\{\lambda^2\unsure{t|k}\mid \lambda\in Z'_n\}$, on est encore loin du compte ! Mais

% archive sheet number: 53
\scanpage{55}

\deleted{notons que pour dans le cas qui nous occupe,}

Mais revenons sur la définition de $Z$ tantôt, en définissant un \inserted{autre} groupe, plus petit encore a priori

\[
 Z_0=\left\{\lambda\in\unsure{\mathfrak S|\mathscr S}_S\ \middle|\
 \begin{aligned}
 &\lambda\text{ centralise }G_0=\operatorname{Aut}(S,\Sigma_0),\\
 &\exists f_\lambda\in\mathfrak S_S\text{ tel qu'on ait, }
 \deleted{\unsure{\forall\{s,t\}\in\mathscr P_2(S)}}\inserted{\forall(s,t)\in S\times S,}\\
 &\{s,t\}\in A_0\Longleftrightarrow
 \{\underbrace{f_\lambda(\deleted{\unsure{x}}s)}_{s'},\underbrace{\lambda^{-1}f_\lambda(\deleted{\unsure{y}}t)}_{t'}\}\in A_0
 \end{aligned}\right.
 \tag*{(138)}
\]

en d'autres termes, $\lambda\in Z_0$ \unsure{sss|ssi} \deleted{pour} \inserted{$\lambda$ centralise $G_0=\operatorname{Aut}(S,\Sigma_0)$, et si pour} $s',t'\in S$, \deleted{la relation} on a

\[
 \bigl(\{s',\lambda^{-1}t'\}\in A_0\Longleftrightarrow
 \{\lambda^{-1}s',t'\}\in A_0\bigr)\Longrightarrow(s'\ne t')
 \tag*{(139)}
\]

de sorte qu'on peut \deleted{désigner par $A$} \inserted{introduire}

\[
 A_\lambda=\left\{\{s',t'\}\in\unsure{\mathscr P|\mathcal P|\mathfrak P}_2(S)\ \middle|\
 \begin{aligned}
 &\{s',\lambda^{-1}t'\}\in A_0\quad\text{i.e.}\\
 &\exists\unsure{(s,t)|\{s,t\}}\in S^2,\ \text{avec }\{s,t\}\in A_0,\\
 &\{s',\deleted{\unsure{s}}t'\}=\{s,\lambda.t\}
 \end{aligned}\right\}
\]

et si $A_\lambda$ est l'ens des arêtes \deleted{pour} une structure triangulée sur $S$ \deleted{\uline{équivalente}} \inserted{isomorphe} à la structure $\Sigma_0$ définie par $A_0\subset\mathscr P_2(S)$, i.e. s'il existe $f_\lambda\in\mathfrak S_S$ telle qu'on ait

\[
 f_\lambda(A_0)=A_\lambda
 \tag*{(140)}
\]

i.e. justement tel que

\[
 \{s,t\}\in A_0\Longleftrightarrow\{f_\lambda(s),\lambda^{-1}f_\lambda(t)\}\in A_0.
 \tag*{(141)}
\]

% archive sheet number: 54
\scanpage{56}
\mspage{27}

. Le fait que $\lambda$ centralise $G_0$ implique que l'on a, pour $u\in G_0$

\[
 \{s',t'\}\in A_\lambda\Longleftrightarrow\{us',ut'\}\in A_\lambda
 \tag*{(142)}
\]

puisque les deux termes de l'équivalence s'écrivent resp.$^{\mathrm t}$

\[
 \{s',\lambda^{-1}t'\}\in A_0\qquad\text{et}\qquad
 \{us',\underbrace{\lambda^{-1}ut^{\unsure{+|\prime}}}_{u\lambda^{-1}t'}\}\in A_0,
\]

relations dont l'équivalence signifie justement que $u\in G_0$. Donc on trouve encore

\[
 \deleted{\unsure{G_\lambda}}\quad\underbrace{\operatorname{Aut}(S,A_\lambda)}_{G_\lambda}=G_0
 \tag*{(143)}
\]

\uline{Définition} Soient $S$ un ensemble, $A_0\subset\unsure{\mathscr P|\mathcal P|\mathfrak P}_2(S)$, $A_1\in\mathscr P_2(S)$ deux structures de graphes \uline{stricts} \inserted{(i.e. une arête connue quand on connaît l'ens. de ses \unsure{extrémités|extrêmités}, et celles-ci sont distinctes)} sur $S$, \inserted{admettant} comme ens. de ses sommets, \inserted{$G_0$ et $G_1$ les groupes des autom. \unsure{correspts.},} On dit que ces structures sont \uline{équivalentes}, s'il existe un $\lambda\in\unsure{\mathfrak S|\mathscr S}_S$, \unsure{telle|tel} que \deleted{l'on} l'on ait

a) $\lambda$ centralise $G_0$

b) \deleted{\unsure{$s,t$}} \deleted{\unsure{Sois}} $s,t\in S$, alors

\[
 \deleted{\unsure{\{\lambda^{-1}s',}}\quad
 \{s,t\}\in A_0\Longleftrightarrow\{s,\lambda t\}\in A_1
\]

% slanted

\margin{On dit alors que $\lambda$ est un \uline{transporteur} de $A_0$ vers $A_1$}

\uline{Remarque} L'argument précédent montre que l'on a alors

\[
 G_0=G_1,
\]

% archive sheet number: 55
\scanpage{57}
% archive sheet number: 56

de sorte que l'on a aussi, pour $s',t'\in S$

\[
 \begin{aligned}
 &\text{a) }\lambda^{-1}\in\text{\unsure{commutateur|commutant} de }G_1\text{ dans }\unsure{\mathfrak S|\mathscr S}_S\\
 &\text{b) }\{s',t'\}\in A_1\Longleftrightarrow\{s',\lambda^{-1}t'\}\in A_0
 \end{aligned}
\]

-- ce qui montre que la relation ``d'équivalence'' introduite est bien symétrique. Elle est évidemment réflexive, et la transitivité se voit de $=$ en utilisant que si \inserted{$A_0\sim A_1$, $A_1\sim A_2$,} \deleted{\unsure{on a}} \deleted{\unsure{$G_0=G_1$}} on aura

\[
 G_0=G_1=G_2
\]

et si

\[
 \left.\begin{aligned}
 \lambda\ &\text{transporte }A_0\text{ en }A_1\\
 \mu\ &\rule[.45ex]{3.1em}{.35pt}\,A_1\text{ en }A_2
 \end{aligned}\right|\quad\lambda\mu\text{ transporte }A_0\unsure{-|\to|\text{en}}A_2
\]

\deleted{Notons que} \inserted{Considérons,} pour $A_0$ fixé, \deleted{le groupe} \inserted{l'ens $\widetilde Z_0$ des} \deleted{des} $\lambda$ tels qu'il \deleted{l'on ait} existe un $A_1$ (évidemment unique) équivalent à $A_0$, et tel que $\lambda$ soit un transporteur de $A_0$ à $A_1$, i.e. \inserted{$\lambda$} tels que

\[
 \left\{\begin{aligned}
 &\lambda\in\operatorname{Comm}_{\mathfrak S_S}(G_0)\\
 &\forall s',t'\in S\quad
 (\{s',\lambda^{-1}t'\}\in A_0)\Longleftrightarrow
 (\{\lambda^{-1}s',t'\}\in A_0)\Longrightarrow s'\ne t'
 \end{aligned}\right.
 \tag*{(144)}
\]

% slanted

\margin{$\forall\unsure{\{s,t\}|s,t}\in S$, on a $\{s,t\}\in A_0\Longleftrightarrow\{\lambda s,\lambda^{-1}t\}\in A_0$ $\Longrightarrow t\ne\lambda s$\quad$\Longleftrightarrow$}

\deleted{est} \inserted{On aimerait que ce soit} un \uline{sous-groupe} de $\mathfrak S_S$, soit $Z_0$.

On \uline{voit} en tous cas que

\[
 \left\{\begin{aligned}
 &\deleted{\unsure{Z_0\cap G_0=\operatorname{Cent}(G_0)\ \text{ab.}}}
 &&\mathrm{id}_S\in\widetilde Z_{\unsure{0|3}}\\
 &\lambda\in\widetilde Z_0\Longrightarrow\lambda^{-1}\in\widetilde Z_0,\qquad
 &&\lambda,\mu\in\unsure{Z_0|\widetilde Z_0},\ [\lambda,\mu]=1\Longrightarrow\lambda\mu\in\widetilde Z_0
 \end{aligned}\right.
 \tag*{(145)}
\]
\scanpage{58}
\mspage{28}
% archive sheet number: 57

donc si \inserted{$\widetilde Z_0$ est formé d'éléments permutables, p.ex. si} $\operatorname{Comm}_{\unsure{\mathfrak S|\mathscr S}_S}(G_0)$ est commutatif, alors $\widetilde Z_0$ est bien un sous-groupe de $\mathfrak S_S$.

% slanted

\margin{\uline{NB} Si $s\mapsto\unsure{\text{Ét}|\text{Et}}_{\Sigma_0}(s)$ $S\longrightarrow\unsure{\mathscr P|\mathcal P|\mathfrak P}(S)$ est injective, alors pour $\lambda,\mu\in\widetilde Z_0$, on a $\lambda\mu\in\widetilde Z_0\Longleftrightarrow\unsure{\mu\lambda}=\lambda\mu$.}

\uline{Définition} On dit que $A_0$ et $A_1$ sont strictement équivalents, s'ils sont équivalents et isomorphes, i.e. s'il existe \deleted{des} $\lambda,\deleted{\unsure{\in}}f\in\mathfrak S_S$, tels que

\[
 \begin{aligned}
 &\text{a) }\lambda\in\operatorname{Comm}_{\mathfrak S_S}(G_0)\\
 &\text{b) }\forall s,t\in S,\text{ on a}\\
 &\{s,t\}\in A_0\Longleftrightarrow\{f(s),f(\unsure{t|M})\}\in A_1\\
 &\phantom{\{s,t\}\in A_0}\Longleftrightarrow\deleted{\unsure{\{}}\{s,\lambda t\}\in A_1
 \end{aligned}
\]

\deleted{\unsure{La}} C'est encore une relation d'équivalence.

Soit $Z_0$ \deleted{\unsure{$\subset\widetilde Z_0$}} l'ens des \deleted{$\lambda\in\mathfrak S_S$} \inserted{$\lambda\in\widetilde Z_0$} tels que \deleted{il} \deleted{existe} le $A_1$ qu'il définit soit \uline{strictement} équivalent à $A_0$, i.e. soit isom. à $A_0$.

sous réserve que $Z_0$ soit formé d'élément \unsure{permutable|permutables}, $Z_0$ est un sous-groupe de $\mathfrak S_S$. On \unsure{suppose|supposera} dorénavant qu'il en est ainsi ! (i.e. que $Z_0$ est formé d'éléments permutables). \deleted{Alors l'ens}

% slanted

\margin{\uline{NB} Si $s\mapsto\text{Ét}(s)$ $S\longrightarrow\mathscr P(S)$ est injective, alors $\unsure{\mathfrak z|z}=$ ens des éléments involutifs de $\operatorname{Cent}(G_0)$}

Soit

\[
 \mathfrak z_0=Z_0\cap G_0=
 \left\{u\in G_0\ \middle|\ \left\{
 \begin{aligned}
 &\inserted{u\in\operatorname{Cent}(G_0)}\\
 &\deleted{\unsure{\{s,t\}}}\ \text{Ét}_{\Sigma_0}(s)=\text{Ét}_{\Sigma_0}(u^2s)\\
 &\forall s\in S
 \end{aligned}\right.\right\}
 \tag*{(146)}
\]
\scanpage{59}

alors l'ens des structures \inserted{\unsure{strict$^{\mathrm{mt}}$|strict$^{\mathrm t}$}} équivalentes à $A_0$ est un torseur sous \deleted{\unsure{$Z_0/\mathfrak z=Z$|$Z_0/\mathfrak z_0$}}

\[
 Z'_0\overset{\mathrm{def}}{=}Z_0/\unsure{\mathfrak z|z}_0
 \tag*{(147)}
\]

Donnons nous, de façon plus intrinsèque, une classe d'équivalence stricte \inserted{$\Pi$} de structures de graphes sur $S$, ``admissibles'' en ce sens que \deleted{la \unsure{com.} sous groupe de $\unsure{\mathfrak S|\mathscr S}_S$} \inserted{l'ens des transporteurs} d'une \deleted{\unsure{structure de}} $A_0\in\Pi$ vers \deleted{des} autres \deleted{\unsure{autres} $A_1\in\Pi$ soit} $A_1\in\Pi$ soit un groupe, soit $Z_0\subset\mathfrak S_S$. Ce groupe ne dépend pas du choix de $A_0$, \unsure{pas|plus} $=$ que $G_0=\operatorname{Aut}(S,A_0)$ ne dépend pas du choix de $A_0\in\Pi$.

% slanted

\margin{$\Pi\in\unsure{\mathscr P|\mathcal P|\mathfrak P}(\mathscr P(\mathscr P_2(S)))$, i.e. $\Pi\subset\underbrace{\mathscr P(\mathscr P_2(S))}_{\text{ens des structures de graphes}}$}

Soit alors

\[
 \begin{aligned}
 \deleted{\unsure{G_f}}\unsure{\mathscr G|\mathcal G}&=\operatorname{Aut}(S,\Pi)
 \overset{\mathrm{def}}{=}
 \{f\in\mathfrak S_S\mid
 \deleted{\forall A_0\in\Pi,\ f(A_0)\in\Pi\text{ et }f^{-1}(A_0)\in\Pi}
 \inserted{f(\Pi)=\Pi}\}\\
 &=\left\{f\in\mathfrak S_S\ \middle|\ 
 \begin{aligned}
 &\forall A_0\in\Pi,\ f(\unsure{F_0|A_0})\in\Pi\\
 & f^{-1}(A_0)\in\Pi
 \end{aligned}\right\}
 \end{aligned}
 \tag*{(148)}
\]

En termes d'\uline{un} $A_0\in\Pi$ choisi, on a

\[
 \mathscr G=\left\{f\in\mathfrak S_S\ \middle|\
 \begin{aligned} &f(A_0)\in\Pi\\ &f(Z_0)=Z_0\end{aligned}\right\}.
 \tag*{(149)}
\]

% archive sheet number: 58
\scanpage{60}
\mspage{29}

On \deleted{\unsure{pourra|trouve}} s'intéressera au ss-groupe $G$ de $\unsure{\mathscr G|\mathcal G}$ défini par

\[
 \deleted{\unsure{\mathscr G}}G=\{f\in\mathscr G\mid f\text{ centralise }Z_0
 \inserted{\bmod\unsure{\mathfrak z|z}_0}\ \text{i.e. }\
 \operatorname{int}(f)|Z_0/\mathfrak z_0=\operatorname{id}_{Z_0/\mathfrak z_0}\}
 \tag*{(150)}
\]

Donc on trouve une suite exacte canonique

\diagram{\begin{tikzcd}[ampersand replacement=\&,column sep=1.15em,row sep=1.8em]
{(151)}\&1\arrow[r]\&G_0\arrow[r]\&G\arrow[r]\&Z'_0\arrow[r]\&1\\
\&1\arrow[r,-]\&\mathfrak z_0\arrow[u,arrows={Hooks[harpoon]-},"{\text{cent.}}"']\arrow[r]\&
 Z_0\arrow[u,hook]\arrow[r,-]\& Z_0/\mathfrak z_0\arrow[u,equal]\arrow[r,-]\&1\arrow[u,phantom,"{\unsure{\wr|=}}" description]{,}
\end{tikzcd}}

\deleted{sans que je sois sûr si $G\to Z'_0$ est surjectif} indépendante du choix d'une ``origine'' $A_0\in\Pi$.

Supposons maintenant que le groupe $G_0$ soit transitif sur l'ens \inserted{$A_0$} des arêtes d'une des structures $A_0\in\Pi$, donc de $=$ \inserted{pour} toutes. On voit alors que

\[
 A_0,A_1\in\Pi,\quad A_0\ne A_1\Longrightarrow
 \underbrace{A_0\cap A_1=\emptyset}_{\substack{\text{intersection de}\\
 \text{parties de }\unsure{\mathscr P|\mathcal P|\mathfrak P}_2(S)}}
 \tag*{(152)}
\]

donc \deleted{si} on pose

\[
 A\unsure{:|=} A_\Pi=\bigcup_{\unsure{A'|A_i}\in\Pi}A'\subset\mathscr P_2(S)
 \tag*{(152)}
\]

on a une réunion disjointe donc une application

\[
 A\xrightarrow{\unsure{\tau|t}}\underset{\text{un }Z'_0\text{-torseur}}{\Pi}
 \tag*{(153)}
\]

% archive sheet number: 59
\scanpage{61}

qui généralise (126).

\deleted{Quand on part d'une situation \unsure{isomorphe} à un}

Pour appliquer ces réflexions au cas d'un $M$ \deleted{\unsure{libre sur|libre}} de rang 2 sur un anneau $\Lambda_n=\mathbb Z/n\mathbb Z$, il faudrait vérifier que pour

\[
 i=\{e,-e\}\in(\text{dét}\,M)^{*}/\{\pm1_\Lambda\}
\]

d'où un $A_{M,i}$, \deleted{\unsure{les élts de son}} \inserted{$\simeq\deleted{\unsure{R}}\operatorname{Rep}(M)/\pm1\subset\unsure{\mathscr P|\mathcal P|\mathfrak P}_2(S_{\unsure{M}})$} où $S=S_M=M^{*}/\{\pm\mathrm{id}_M\}$, on a \deleted{déjà}

\[
 \underset{\substack{\cap\\ \unsure{\mathfrak S|\mathscr S}_S}}{Z_0}
 =\deleted{\unsure{Z'_n\ \text{opérant par}}}
 \underbrace{\{T_\lambda\mid\lambda\in{Z'_n}=\Lambda^*/\pm1_\Lambda\}}_{\text{noté }(Z'_n)_S},
 \tag*{(154)}
\]

et on suspecte même que l'on a déjà

\[
 \operatorname{Comm}_{\mathfrak S_S}
 \bigl(\underset{\substack{\Vert\\{\mathrm{Gl}}^{\unsure{1|\prime}}_\Lambda(M)^\pm}}{G_0}\bigr)
 \overset{?}{=}(Z'_n)_S\ \unsure{;|.}
 \tag*{(155)}
\]

Mais sans s'encombrer de l'hypothèse $\Lambda\simeq\mathbb Z/n\mathbb Z$, on peut regarder sur $S_M=\underbrace{M^{**}}_{\substack{\text{ens des}\\\text{él. de }M\\\text{qui font partie d'une}\\\text{base}}}/\{\pm\mathrm{id}_M\}$ la structure de graphe définie par $A_{M,i}$ $(i=\{e,-e\}\in(\text{det}\,M)^*/\pm1_\Lambda))$ et les groupes $G_0,Z_0$ associés.

% archive sheet number: 60
\scanpage{62}
\mspage{30}

$\underline{\overline{\mathrm{IV}}}$ Soient \deleted{\unsure{$R$}} $\Lambda$ un anneau non nul, $M$ un module libre de rang 2 sur $\Lambda$,

\[
 \left\{\begin{aligned}
 \widetilde S_M&=\{x\in M\mid\exists y\in M,\ (x,y)\in\operatorname{Rep}(M)\}\\
 \deleted{\text{(157)}}\qquad S_M&=\widetilde S_M/(\pm\mathrm{id}_M)
 \end{aligned}\right.
 \tag*{(156)}
\]

de telle façon que l'hom canonique

\[
 \left\{\begin{aligned}
 &\widetilde S_M\longrightarrow S_M\\
 &\text{est de degré constant }\left\{\begin{aligned}
 &\text{égal à }2\text{ si }2.1_\Lambda\ne0\\
 &\text{égal à }1\text{ si }2.1_\Lambda=0,\ \text{i.e.}\\
 &\Lambda\text{ est une }\unsure{\mathbb F_2|\mathrm F_2}\text{-alg.}
 \end{aligned}\right.
 \end{aligned}\right.
 \tag*{(157)}
\]

\unsure{Soient|Soit}

\[
 \left\{\begin{aligned}
 &\deleted{\widetilde A_M=\operatorname{Rep}(M)\ \unsure{\operatorname{Rep}(M)}/\text{symétrie }(x,y)\mapsto(y,x)}\\
 &\deleted{\phantom{\widetilde A_M}=\unsure{R}\{\{x,y\}\in\unsure{\mathscr P|\mathcal P|\mathfrak P}_2(M)\mid\{x,y\}\text{ une base de }M\}}\\
 &A_M=\deleted{\operatorname{Rep}(M)/\unsure{\pm\mathrm{id}_{\pm1}|\pm\mathrm{id}_M}\ \widetilde A_M/\{\pm\mathrm{id}_M\}}\\
 &\quad\left\{\{s,t\}\in\mathscr P_2(S_M)\ \middle|\ 
 \begin{aligned}&\exists(x,y)\in\operatorname{Rep}(M),\\ &\inserted{\{s,t\}=\{\dot x,\dot y\}}\end{aligned}\right\}
 \end{aligned}\right.
 \tag*{(158)}
\]

\margin{\uline{NB} Il y a \textcircled{4} \uline{choix} (et non 2 !) de \deleted{\unsure{une telle}} \inserted{$x,y$} tels que $(s,t)=(\dot x,\dot y)$, \inserted{donc 8 choix} \deleted{tels} que $\{s,t\}=\{\dot x,\dot y\}$}

\deleted{de sorte que l'on a}

\[
 \deleted{\left\{\begin{smallmatrix}
 \widetilde A_M\unsure{-|\to|\longrightarrow} A_M\\
 \text{de degré constant }\left\{\begin{smallmatrix}
 \text{égal à }2\text{ si }2.1_\Lambda\ne0\\
 \rule[.45ex]{2.3em}{.35pt}\,1\text{ si }2.1_\Lambda=0.
 \end{smallmatrix}\right.
 \end{smallmatrix}\right.}
 \tag*{\deleted{(159)}}
\]

\deleted{On a une inclusion canonique}

\[
 \deleted{\overset{\deleted{\unsure{\sim}}}{A}_M\lhook\joinrel\longrightarrow\mathscr P_2(S_M)}
 \tag*{\deleted{(\unsure{160|150})}}
\]

\deleted{du fait que si $\{x,y\}$ est une base de $M$, alors $y\notin\deleted{\unsure{x}}\{x,-x\}$ donc $\{\dot x,\dot y\}\in\mathscr P_2(S_M)$, \deleted{et bien sûr $\{\dot x,\dot y\}\unsure{=\{x,|=\{-x,}$,} soit $\{x,y\}^{\unsure{\bullet|\cdot}}$, et bien sûr on a}

\[
 \deleted{\bigl(\underbrace{-\{x,y\}}_{\overset{\mathrm{def}}{=}\{-x,-y\}}\bigr)^{\cdot}
 =\{x,y{\unsure{\}|\rbrack}}.}
\]

% archive sheet number: 61
\scanpage{63}

On considère $A_M$ comme l'ens des \uline{arêtes} d'une structure de graphe sur $S_M$. \deleted{Le groupe} On a une application naturelle

\[
 A_M \xrightarrow{\unsure{\tau|t}\text{ ou }\tau_A} \deleted{\unsure{\Delta|A}}\Pi_M \overset{\text{déf}}{=} (\text{dét}\,M)^*/{\pm 1_\Lambda}
 \tag*{(161)}
\]

% slanted
\margin{on regardera $\Pi_M$ comme un ens. quotient de $A_M$, \unsure{grâce} à $\tau$.}

déduite par passage au\deleted{\unsure{x}} quotient de l'application

\[
 \left\{\begin{aligned}
  \widetilde A_M &\xrightarrow{\widetilde\tau} \Pi_M \\
  \{x,y\} &\longmapsto x\wedge y \bmod\text{signe}
 \end{aligned}\right.
 \tag*{(162)}
\]

\margin{NB. ne change pas si on échange $x$ et $y$ ou si on change $x$ ou $y$ ou les deux de signe}

Posons, pour $i\in\Pi_M$

\[
 A_{M,i}=\tau^{-1}(\{i\})\ , \tag*{(163)}
\]

qu'on considère comme l'ens des arêtes pour une structure de graphe \deleted{\unsure{sur}} $\Sigma(i)$ sur $S_M$.

Le groupe $\operatorname{Gl}_\Lambda(M)$ opère sur $(\deleted{\unsure{A}}\unsure{S|A}_M,A_M)$ par transport de structure, ainsi que sur $\Pi_M$, et $\tau$ (161) commute aux opérations de $\operatorname{Gl}_{\unsure{\Lambda|\lambda}}(M)$. Comme le sous-groupe $\{\pm\operatorname{id}_M\}$ de $\operatorname{Gl}_\Lambda(M)$ opère trivialement sur les ens $S_M,\Pi_M$, on

% archive sheet number: 62
\scanpage{64}
\mspage{31}

conclut que $\operatorname{Gl}_\Lambda(M)$ opère par l'intermédiaire d'une \deleted{opération}

\[
 \begin{aligned}
 &\text{opération de } \inserted{G_M^{\deleted{\unsure{*}}}=}\operatorname{Gl}_\Lambda(M)' \inserted{=\operatorname{Gl}_\Lambda(M)/{\pm1}}\text{ sur } S_M,\\
 &\text{respectant la structure } A_M\subset\unsure{\mathscr P|\mathcal P|\mathfrak P}_2(S_M),\\
 &\text{et l'ens. quotient }\Pi_M\text{ de }A_M,
 \end{aligned}\tag*{(164)}
\]

% slanted
\margin{\uline{NB} l'opération de $G_M$ sur $\Pi_M$ respecte aussi sa structure de $Z'_\Lambda$-torseur, \deleted{\unsure{où}} $Z'_\Lambda=\Lambda^*/{\pm1_\Lambda}$, et $G_M$ opère sur $\Pi_M$ via $\delta_M$.}

pour $i\in\Pi$, $g\in G_M$, $g_{S_M}$ est un isom. de la structure $(S_M,A_{M,i})$ sur la structure $(S_M,A_{M,g_{\Pi_M}i})$. En particulier, si on \deleted{\unsure{pose}} considère le composé

\diagram{\begin{tikzcd}[ampersand replacement=\&]
 \text{(165)} \& G_M \arrow[r,"\text{det}"] \arrow[rr,bend right=15,"\delta_M"'] \& \Lambda^* \arrow[r,"\text{can}"] \& \Lambda^*/{\pm1_\Lambda}\overset{\text{def}}{=}Z'_\Lambda{,}
\end{tikzcd}}

le stabilisateur \inserted{dans $G_M$} de \deleted{tout} point de $\Pi_M$ \deleted{est} le noyau \inserted{$G_M^0={}$} $\operatorname{Ker}\delta_M$ intervenant dans la suite exacte

\diagram{\begin{tikzcd}[ampersand replacement=\&]
 \text{(\unsure{166|165})} \& 1\arrow[r,-] \& G_M^0\arrow[r] \& G_M\arrow[r,"\delta_M"] \& Z'_\Lambda\arrow[r,-] \& 1{,}\\
 \& \& {\unsure{\operatorname{Sl}|\operatorname{Gl}}}_\Lambda^1(M)^{\unsure{\pm1|/{\pm1}}}\arrow[u,equal] \& \operatorname{Gl}_\Lambda^{\unsure{*}}(M)'\arrow[u,equal]
\end{tikzcd}}

et il en résulte que c'est aussi le sous-groupe de $G_M$ qui stabilise la structure de \deleted{graphe} graphe $\Sigma(i)$, pour un $i\in\Pi_M$ quelconque donné.

% archive sheet number: 63
